\documentclass[12pt,a4paper]{article}
\usepackage[utf8]{inputenc}
\usepackage{amsmath,amssymb,amsthm}
\usepackage{algorithm}
\usepackage{algpseudocode}
\usepackage{geometry}
\geometry{margin=1in}

\newtheorem{theorem}{Theorem}
\newtheorem{lemma}[theorem]{Lemma}

\title{Problem 66: Diophantine Equation (Pell's Equation)}
\author{Project Euler Solution}
\date{}

\begin{document}
\maketitle

\section{Problem Statement}

Consider quadratic Diophantine equations of the form:
\[
x^2 - Dy^2 = 1
\]
Find the value of $D \le 1000$ in minimal solutions of $x$ for which the largest value of $x$ is obtained.

\section{Theory: Pell's Equation and Continued Fractions}

\subsection{Periodic Continued Fractions}

\begin{theorem}[Lagrange]
For any non-square positive integer $D$, the continued fraction expansion of $\sqrt{D}$ is eventually periodic:
\[
\sqrt{D} = [a_0; \overline{a_1, a_2, \ldots, a_{r-1}, 2a_0}]
\]
where $a_0 = \lfloor \sqrt{D} \rfloor$ and $r$ is the period length.
\end{theorem}

\subsection{Convergents}

The $k$-th convergent $p_k / q_k$ of the continued fraction is computed via:
\begin{align}
p_{-1} &= 1, \quad p_{-2} = 0, \quad p_k = a_k p_{k-1} + p_{k-2} \\
q_{-1} &= 0, \quad q_{-2} = 1, \quad q_k = a_k q_{k-1} + q_{k-2}
\end{align}

\subsection{Fundamental Solution}

\begin{theorem}
Let $r$ be the period of the continued fraction expansion of $\sqrt{D}$. The fundamental (minimal positive) solution $(x_1, y_1)$ of $x^2 - Dy^2 = 1$ is:
\[
(x_1, y_1) = \begin{cases}
(p_{r-1}, q_{r-1}) & \text{if } r \text{ is even} \\
(p_{2r-1}, q_{2r-1}) & \text{if } r \text{ is odd}
\end{cases}
\]
\end{theorem}

\begin{proof}
The convergents of $\sqrt{D}$ satisfy $p_k^2 - D q_k^2 = (-1)^{k+1} \cdot c_k$ where $c_k$ divides $2a_0$. At the end of the period ($k = r-1$), we have $p_{r-1}^2 - D q_{r-1}^2 = (-1)^r$.

If $r$ is even, then $p_{r-1}^2 - D q_{r-1}^2 = 1$, giving the fundamental solution directly.

If $r$ is odd, then $p_{r-1}^2 - D q_{r-1}^2 = -1$. We use the identity:
\[
(p_{r-1}^2 + D q_{r-1}^2)^2 - D(2 p_{r-1} q_{r-1})^2 = (p_{r-1}^2 - D q_{r-1}^2)^2 = 1
\]
This corresponds to the convergent at index $2r-1$, which can be verified by the periodicity of the continued fraction coefficients.
\end{proof}

\subsection{Continued Fraction Algorithm}

To compute the continued fraction of $\sqrt{D}$, we use:
\begin{align}
m_0 &= 0, \quad d_0 = 1, \quad a_0 = \lfloor \sqrt{D} \rfloor \\
m_{n+1} &= d_n a_n - m_n \\
d_{n+1} &= \frac{D - m_{n+1}^2}{d_n} \\
a_{n+1} &= \left\lfloor \frac{a_0 + m_{n+1}}{d_{n+1}} \right\rfloor
\end{align}
The algorithm terminates when $a_n = 2a_0$ (end of period).

\section{Algorithm}

\begin{algorithm}[H]
\caption{Find $D \le 1000$ with largest minimal $x$}
\begin{algorithmic}[1]
\State $x_{\max} \gets 0$, $D_{\text{best}} \gets 0$
\For{$D = 2$ to $1000$}
    \If{$D$ is a perfect square} \textbf{continue} \EndIf
    \State Compute continued fraction of $\sqrt{D}$
    \State Compute convergents until $p_k^2 - D q_k^2 = 1$
    \If{$p_k > x_{\max}$}
        \State $x_{\max} \gets p_k$, $D_{\text{best}} \gets D$
    \EndIf
\EndFor
\State \Return $D_{\text{best}}$
\end{algorithmic}
\end{algorithm}

\section{Complexity Analysis}

\begin{itemize}
    \item \textbf{Time}: $O(N \cdot \sqrt{N} \log N)$ where $N = 1000$, since the period length of the continued fraction for $\sqrt{D}$ is $O(\sqrt{D} \log D)$.
    \item \textbf{Space}: $O(1)$ auxiliary per $D$, since we only track two consecutive convergents.
    \item We use arbitrary-precision integers (in Python/C++ with big integers) since $x$ can be extremely large (up to $\sim 10^{37}$).
\end{itemize}

\section{Result}

The answer is $\boxed{D = 661}$.

The minimal $x$ for $D = 661$ is:
\[
x = 16421658242965910275055840472270471049
\]

\section{Answer}

\[
\boxed{661}
\]

\end{document}
